\documentclass{ximera}
\title{The Product Rule}

\newcommand{\pskip}{\vskip 0.1 in}

\begin{document}
\begin{abstract}
An introduction to the product rule.
\end{abstract}
\maketitle


\section{The Product Rule, a Special Case}

\begin{question} \label{Q23mcm32}
The purpose of this question is to prove the relative product rule in a special case. Do \emph{not} assume it to be true.

The functions $P=f(t)$ and $Q=g(t)$ change at constant relative rates. 

\begin{enumerate}

\item Suppose $f(0)=P_0$ and that
\[
   \frac{1}{P} \cdot \frac{dP}{dt} = a. 
\]
Find an expression for the function $P=f(t)$.
\[
    P = f(t) = P_0 \answer{e^{at}} .
\]

\item Suppose $g(0)=Q_0$ and that
\[
   \frac{1}{Q} \cdot \frac{dQ}{dt} = b. 
\]
Find an expression for the function $Q=g(t)$.
\[
    P = f(t) = Q_0 \answer{e^{bt}} .
\]

\item Find an expression for the function $R = PQ = f(t)g(t)$.
\[
      R = PQ = P_0Q_0e^{(\answer{a+b})t} .
\]

\item Express the relative instantaneous rate of change of the product 
\[
  R=PQ = f(t)g(t)
\]
in terms of the relative instantaneous rates of change of its factors $P$ and $Q$. Explain your reasoning.

\item State a clear conclusion without mathematical notation. 

If two functions change at constant relative rates, then the relative instantaneous rate of change of their product is equal to ....

\begin{freeResponse}
\end{freeResponse}
\end{enumerate}

\end{question}

\section{The Natural Log Function and its Derivative}

\begin{question} \label{Qf0e0feef0}
\begin{enumerate}
\item Use the methods of this class, not implicit differentiation, to find the derivative of the inverse of the function
\[
     P = f(t) = e^t \, ,\, t\in \mathbb{R}.
\]
Include the correct domain.

Start by expressing the derivative $dP/dt$ in terms of $P$.

\item Find simplified expressions for each of the following derivatives. Do each problem twice. First directly as is, then again by simplifying each function before differentiating.

\begin{enumerate}
\item $\frac{d}{dx} \left( \ln \left( \frac{x}{4} \right) \right)$

\item $\frac{d}{dx} \left( \ln x^4 \right)$

\item $\frac{d}{d\theta} \left(\ln \left| \sec\theta \right|  \right)$
\end{enumerate}
\end{enumerate}
\end{question}


\begin{question} \label{Qp0v3f30}
Let  $P=f(t)>0$ be a differentiable function of $t$.  Use the Leibniz notation to find an expression for the derivative
\[
         \frac{d}{dt} \left( \ln P \right) =   \frac{d}{dt} \left( \ln (f(t)) \right) 
\]
and interpret its meaning. 
\end{question}


\begin{question} \label{Qppv3frr}
Suppose $P=f(t)$ and $Q=g(t)$ are positive differentiable functions of $t$. 

\begin{enumerate}
\item Use the result of Question 3 to express the relative rate of change of the product 
\[
   R = h(t) =  PQ = f(t)g(t)
\]
in terms of the relative rates of change of $P$ and $Q$.

\item Use the result of part (a) to express the rate of change of the product $R=PQ$ in terms of $P$, $Q$, and their rates of changes. Check your expression has the correct units.

\item State a clear conclusion without mathematical notation. What can you say about the relative instantaneous rate of change of the product of two differentiable functions?

\begin{freeResponse}
\end{freeResponse}

\end{enumerate}
\end{question}


\section{Demand Functions and the Product Rule}

\begin{question} \label{Qpvm3f33frnb9}
The function 
\[
    Q = f(P) \, , \, 6 \leq P \leq 15,
\]
expresses the average number of burgers Five Guys of Edmonds sells (measured in burgers/day) in terms of the price (measure in dollars/burger).

Suppose 
\[
     f(10) = 200
\]
and
\[
   \frac{dQ}{dP}\Big|_{P=10} = -16.
\]

\begin{enumerate}
\item Interpret the meaning of the derivative above as a scaling factor. Give an example with \emph{specific} changes.

\item Evaluate the derivative 
\[
      \frac{d}{dP} \left( \ln Q  \right)\Big|_{P=10}
\]
and interpret its meaning as a scaling factor. Give an example with \emph{specific} changes.

\item Find the scaling factor that converts (by multiplication) a small relative change in the price (measured in dollars/burger, from the price of $\$10$/burger) to an accurate approximation of the corresponding relative change in the average number of burgers sold/day. Include units. Give an example with \emph{specific} changes.

\end{enumerate}
\end{question}

\begin{question} \label{Qefefrf22}
This is a continuation of the previous. Use the information given there.

Let 
\[
   R = g(P) \, , \, 6\leq P \leq 15,
\]
be the function that expresses the average revenue (measued in dollars/day) from the sale of burgers in terms of the price (measured in dollars/burger).

\begin{enumerate}
\item Use numerical methods to approximate the derivative
\[
 \frac{dR}{dP}\Big|_{P=10}.
\]
Include \emph{unsimplified} units. Do this by filling in the missing entries in the table below.

Note the entries in the second, third, and fourth columns will be approximations.

\begin{center}
  \begin{tabular}{ | c| c | c | c | c |}
    \hline
    $P$ (dollars/burger) & $Q$ (burgers/day) & $R=PQ$ (dollars/day) & $\Delta R/\Delta P$ (units?) \\ \hline
      $9$ & $\answer{216}$ & $\answer{1944}$ &  $\answer{56}$ \\ \hline
    $9.9$ & $\answer{201.6}$ & $\answer{1995.84}$ &  $\answer{41.6}$  \\ \hline
    $9.99$  & $\answer{200.16}$  & $\answer{1999.5984}$  & $\answer{40.16}$   \\ \hline
           \hline
  \end{tabular}
\end{center}

\item Use calculus to evaluate the derivative
\[
       \frac{dR}{dP} \Big|_{P=10}
\]
and interpret its meaning in the usual way.


\item What are the units of the derivative in part (b)? Explain its meaning in the usual way. Do \emph{not} simplify the units.

\item Evaluate the derivative
\[
       \frac{d}{dP} \left( \ln P \right) \Big|_{P=10}
\]
and interpret its meaning in the usual way. 

\item Evaluate the derivative
\[
       \frac{d}{dP} \left( \ln R \right) \Big|_{P=10}
\]
and interpret its meaning in the usual way.


\item If you were the manager and Five Guys were currently charging $\$10$/burger, would you raise the price? Explain.
\begin{freeResponse}
\end{freeResponse}

\item Simplify the units of the derivative in part (a). What, if anything, does this suggest about its meaning?
\begin{freeResponse}
\end{freeResponse}
\end{enumerate}

\end{question}

\section{Back to the Quotient Rule}
\begin{question} \label{Qppf3r39c}
Let $P=f(t)$ and $Q=g(t)>0$ be differentiable functions of $t$.

Let 
\[
  R = h(t) = \frac{f(t)}{g(t)} = \frac{P}{Q}.
\]

Use the algebra of limits to find an expression for the derivative
\[
 \frac{d}{dt} \left( \frac{f(t)}{g(t)} \right) = \frac{d}{dt} \left( \frac{P}{Q} \right) .
\]

Do this as follows:

\begin{align*}
 \frac{d}{dt} \left( \frac{P}{Q} \right)   &= \lim_{\Delta t \to 0} \frac{\Delta (P/Q)}{\Delta t} \\
                                                        &= \lim_{\Delta t \to 0} \frac{1}{\Delta t} \cdot \left( \frac{P+\Delta P}{Q+\Delta Q} - \frac{P}{Q}   \right) \\
                                                       &= \lim_{\Delta t \to 0} \frac{1}{\Delta t} \cdot \left( \frac{\answer{Q\Delta P - P\Delta Q}}{Q(Q+\Delta Q)} \right) \\
                                                     &=\lim_{\Delta t \to 0} \frac{\Delta P}{\Delta t}\cdot \frac{1}{\answer{Q+\Delta Q}} - \lim_{\Delta t \to 0}\frac{P}{\answer{Q(Q+\Delta Q)}} \cdot \frac{\Delta Q}{\Delta t} \\
                                                     &= \frac{dP}{dt}\cdot \frac{1}{\answer{Q}} - \frac{P}{\answer{Q^2}}\cdot \frac{dQ}{dt} .
\end{align*}

\end{question}



\end{document}